\chapter{Appendix}\label{app:Appendix}
\begin{longtblr}[caption = {Full detailed table of the dataset}]{
    colspec = {|m{0.5\linewidth} m{0.15\linewidth}  m{0.12\linewidth} m{0.12\linewidth}|},
    rowhead = 1}
        \hline
        \textbf{Class} & \textbf{Type} & \textbf{Total number of samples} & \textbf{Number of sample in testset} \\ \hline
        Gryllteiste & bird & 4 & 1 \\ \hline
        Offshore Windpark & anthro & 4 & 1 \\ \hline
        Schnatterente & bird & 4 & 1 \\ \hline
        Buchfink & bird & 4 & 1 \\ \hline
        unbestimmte Larusmöwe & bird & 4 & 1 \\ \hline
        Fischereigerät (Schwimmer) & anthro & 4 & 1 \\ \hline
        Schmarotzer/Spatel/Falkenraubmöwe & bird & 5 & 1 \\ \hline
        Brandgans & bird & 5 & 1 \\ \hline
        Wasserlinie mit Großalgen & misc & 5 & 1 \\ \hline
        unbestimmte Art & bird & 5 & 1 \\ \hline
        Feldlerche & bird & 5 & 1 \\ \hline
        Schmarotzerraubmöwe & bird & 5 & 1 \\ \hline
        Grosser Brachvogel & bird & 6 & 2 \\ \hline
        Öl & misc & 6 & 2 \\ \hline
        unbestimmte Raubmöwe & bird & 7 & 2 \\ \hline
        Turmfalke & bird & 7 & 2 \\ \hline
        Trauerseeschwalbe & bird & 7 & 2 \\ \hline
        Front & misc & 8 & 2 \\ \hline
        unbestimmter Schwan & bird & 8 & 2 \\ \hline
        Boot (unbestimmtes kleines Boot) & anthro & 8 & 2 \\ \hline
        Tonne & anthro & 8 & 2 \\ \hline
        Sperber & bird & 8 & 2 \\ \hline
        Kiebitzregenpfeifer & bird & 8 & 2 \\ \hline
        Skua & bird & 9 & 2 \\ \hline
        Graugans & bird & 9 & 2 \\ \hline
        Freizeitfahrzeug & anthro & 10 & 2 \\ \hline
        unbestimmte Krähe & bird & 10 & 2 \\ \hline
        Pfuhlschnepfe & bird & 11 & 3 \\ \hline
        Frachter & anthro & 11 & 3 \\ \hline
        Ankerboje & anthro & 12 & 3 \\ \hline
        Fischereifahrzeug & anthro & 12 & 3 \\ \hline
        Mantel-/Heringsmöwe & bird & 13 & 3 \\ \hline
        unbestimmter kleiner Wal & mammal & 13 & 3 \\ \hline
        Papageitaucher & bird & 13 & 3 \\ \hline
        unbestimmte Limikole & bird & 14 & 3 \\ \hline
        Versorgungsschiff & anthro & 14 & 3 \\ \hline
        Gänsesäger & bird & 14 & 3 \\ \hline
        Trübungsfahne & misc & 15 & 3 \\ \hline
        Stellnetzfahne rot & anthro & 15 & 3 \\ \hline
        Pfeifente & bird & 15 & 3 \\ \hline
        Fisch: sichtbares Fischvorkommen & misc & 17 & 4 \\ \hline
        Ringelgans & bird & 17 & 4 \\ \hline
        Fahne zur markierung von Fischereifanggeräten & anthro & 19 & 4 \\ \hline
        Mikroalgen (Phytoplankton): Algenwolke unter Wasser & misc & 19 & 4 \\ \hline
        Rothalstaucher & bird & 20 & 4 \\ \hline
        Fahrspur eines Wasserfahrzeugs & misc & 21 & 5 \\ \hline
        Schiff (unbestimmtes großes Schiff) & anthro & 21 & 5 \\ \hline
        Kegelrobbe & mammal & 23 & 5 \\ \hline
        Eiderente & bird & 23 & 5 \\ \hline
        Haubentaucher & bird & 25 & 5 \\ \hline
        Tang / Makroalgen & misc & 25 & 5 \\ \hline
        Mikroalgen: Algenblütenpatch an Wasseroberfläche & misc & 26 & 6 \\ \hline
        Wasserlinie mit Schaum & misc & 26 & 6 \\ \hline
        unbestimmter Singvogel & bird & 27 & 6 \\ \hline
        Lachmöwe & bird & 27 & 6 \\ \hline
        Wasserlinie mit Algenblüte & misc & 27 & 6 \\ \hline
        Sturm/Silbermöwe & bird & 27 & 6 \\ \hline
        unbestimmter Meeressäuger & mammal & 28 & 6 \\ \hline
        Windenergieanlage & anthro & 28 & 6 \\ \hline
        Höckerschwan & bird & 28 & 6 \\ \hline
        Samtente & bird & 29 & 6 \\ \hline
        Arbeitsschiff Windpark & anthro & 29 & 6 \\ \hline
        unbestimmte Robbe (Kegelrobbe/Seehund) & mammal & 29 & 6 \\ \hline
        unbestimmte Möwe & bird & 29 & 6 \\ \hline
        Seehund & mammal & 29 & 6 \\ \hline
        Ohrentaucher & bird & 29 & 6 \\ \hline
        Trottellumme & bird & 29 & 6 \\ \hline
        Trauerente & bird & 29 & 6 \\ \hline
        Mantelmöwe & bird & 30 & 6 \\ \hline
        Mittelsäger & bird & 30 & 6 \\ \hline
        Prachttaucher & bird & 30 & 6 \\ \hline
        Sonstiges & misc & 30 & 6 \\ \hline
        Eisente & bird & 30 & 6 \\ \hline
        Trottellumme/Tordalk & bird & 31 & 7 \\ \hline
        Zwergmöwe & bird & 31 & 7 \\ \hline
        Schweinswal & mammal & 31 & 7 \\ \hline
        Dreizehenmöwe & bird & 31 & 7 \\ \hline
        Silbermöwe & bird & 31 & 7 \\ \hline
        unbestimmter Vogel & bird & 31 & 7 \\ \hline
        Eissturmvogel/Möwe & bird & 32 & 7 \\ \hline
        Marine Lebewesen & misc & 32 & 7 \\ \hline
        Heringsmöwe & bird & 32 & 7 \\ \hline
        Messboje & anthro & 32 & 7 \\ \hline
        Basstölpel & bird & 32 & 7 \\ \hline
        unbestimmte Kleinmöwe & bird & 32 & 7 \\ \hline
        Boje oder Bojenkette & anthro & 32 & 7 \\ \hline
        unbestimmter Alk & bird & 32 & 7 \\ \hline
        Stockente & bird & 33 & 7 \\ \hline
        unbestimmter Seetaucher & bird & 33 & 7 \\ \hline
        Seeschwalbe / Kleinmöwe & bird & 33 & 7 \\ \hline
        unbestimmte Großmöwe & bird & 33 & 7 \\ \hline
        unbestimmte Ente & bird & 33 & 7 \\ \hline
        Windenergieanlage im Bau & anthro & 33 & 7 \\ \hline
        Müll & misc & 33 & 7 \\ \hline
        Trauerente/Samtente & bird & 33 & 7 \\ \hline
        Robbe / kleiner Wal & mammal & 33 & 7 \\ \hline
        Sturmmöwe & bird & 33 & 7 \\ \hline
        Wasserlinie mit Müll & misc & 33 & 7 \\ \hline
        Langmuirstreifen & misc & 33 & 7 \\ \hline
        Luftballon & misc & 33 & 7 \\ \hline
        Qualle(n) & misc & 33 & 7 \\ \hline
        Mantel/Heringsmöwe & bird & 33 & 7 \\ \hline
        Tordalk & bird & 33 & 7 \\ \hline
        Sterntaucher & bird & 33 & 7 \\ \hline
        Eissturmvogel & bird & 34 & 7 \\ \hline
        Goldregenpfeifer & bird & 34 & 7 \\ \hline
        unbestimmte Sterna-Seeschwalbe & bird & 34 & 7 \\ \hline
        Weisswangengans & bird & 34 & 7 \\ \hline
        Blässhuhn & bird & 34 & 7 \\ \hline
        Fischereigerät (Doppel-Schwimmer) & anthro & 34 & 7 \\ \hline
        Kormoran & bird & 34 & 7 \\ \hline
        Schellente & bird & 34 & 7 \\ \hline
        unbestimmte Graue Gans & bird & 34 & 7 \\ \hline
        Fluss-/Küstenseeschwalbe & bird & 34 & 7 \\ \hline
        Bergente & bird & 34 & 7 \\ \hline
        unbestimmte Seeschwalbe & bird & 34 & 7 \\ \hline
        Brandseeschwalbe & bird & 34 & 7 \\ \hline
\end{longtblr}
\label{tab:full_dataset}
% \section{Hyperparameter tuning in augmentations}
% \subsubsection{CLAHE}
% The CLAHE augmentation technique has two hyperparameters: the clipping limit and the filter size. The clipping limit sets the maximum threshold for contrast limiting, while the filter size determines the size of the grid used for histogram equalization. To tune these hyperparameters, we used log-scale values ranging from 2 to 32. The hyperparameter tuning results for CLAHE are presented in table \ref{tab:clahe_hp}. 

\begin{table}
    \caption{Results of hyperparmeter tuning for CLAHE.}
    \label{tab:clahe_hp}
    \centering
    \small
    \tabcolsep=0.11cm
    \begin{tabular}{|cc|c|}
    \hline
        Clipping Limit & Filter Size & F1-Score top-3 \\ \hline
        2 & 4 & 0.098 \\ 
        2 & 8 & 0.0965 \\ 
        2 & 16 & 0.0357 \\ 
        2 & 32 & 0.0607 \\ 
        4 & 4 & 0.1111 \\ 
        4 & 8 & 0.0862 \\ 
        4 & 16 & 0.1189 \\ 
        4 & 32 & 0.085 \\ 
        8 & 4 & 0.1113 \\ 
        8 & 8 & 0.0893 \\ 
        8 & 16 & 0.067 \\ 
        8 & 32 & 0.0543 \\
        16 & 4 & \textbf{0.1194} \\ 
        16 & 8 & 0.0765 \\ 
        16 & 16 & 0.0393 \\ 
        16 & 32 & 0.0282 \\ 
        32 & 4 & 0.0536 \\ 
        32 & 8 & 0.09 \\ 
        32 & 16 & 0.0601 \\ 
        32 & 32 & 0.0333 \\ 
    \hline
    \end{tabular}

\end{table}

% Color Jittering transformation has four hyperparameters: brightness, hue, contrast, and saturation. The values of these parameters determine the upper limit of the random magnitude applied to the image. We explored the maximum bound value for all parameters within the range of 0.2 to 0.8. The resulting hyperparameter tuning outcomes are presented in table \ref{tab:colorjittering_hp}.

\begin{table}
\caption{Results of hyperparmeter tuning for  Color Jittering.}
\label{tab:colorjittering_hp}
    \centering
    \small
    \tabcolsep=0.11cm

\begin{tabular}{|cccc|c|}
\hline
Brightness & Hue & Contrast & Saturation & F1-Score top-3 \\
\hline
 0.57 & 0.45 & 0.56 & 0.78 & 0.0125 \\
 0.74 & 0.73 & 0.6 & 0.4 & 0.0250 \\
 0.41 & 0.28 & 0.55 & 0.47 & 0.0388 \\
 0.55 & 0.2 & 0.23 & 0.64 & 0.0418 \\
 0.67 & 0.26 & 0.65 & 0.27 & 0.0507 \\
 0.77 & 0.77 & 0.24 & 0.36 & 0.0518 \\
 0.48 & 0.55 & 0.34 & 0.64 & 0.0519 \\
 0.8 & 0.42 & 0.54 & 0.65 & \textbf{0.0585} \\
\hline
\end{tabular}
\end{table}

% \subsubsection{GridDropout}
% For Griddropout, we can adjust the number of grids and the size of each hole. The hole ratio refers to the size of the hole relative to a fixed grid cell. However, these two hyperparameters are correlated because a smaller number of holes can result in different interpretations of the same hole ratio. Therefore, we fixed the number of grid holes to 100 and only varied the hole ratio during hyperparameter tuning. The resulting hyperparameter tuning outcomes are presented in table \ref{tab:griddropout_hp}.

\begin{table}
\caption{Results of hyperparmeter tuning for GridDropout.}
    \label{tab:griddropout_hp}
    \centering
    \small
    \tabcolsep=0.11cm
    \begin{tabular}{|c|c|}
    \hline
        Hole ratio & F1-Score top-3  \\ \hline
        0.6 & 0.1021  \\ 
        0.5 & \textbf{0.1116}  \\ 
        0.2 & 0.08904  \\ 
        0.3 & 0.1013  \\ 
        0.4 & 0.0914  \\
        \bottomrule
    \end{tabular}
    
\end{table}

% \subsubsection{Sharpen}
% In sharpen we have to determine the alpha range and the lightness range. Alpha is the range to choose the visibility of the sharpened image. At 0, only the original image is visible, at 1.0 only its sharpened version is visible. Lightness is the range to choose the lightness of the sharpened image. We do grid search for same-size range with different of 0.2 from 0.2 to 1.0. The results are presented in table \ref{tab:sharpen_hp}.

\begin{table}
\caption{Results of hyperparmeter tuning for Sharpen.}
    \label{tab:sharpen_hp}
    \centering
    \small
    \tabcolsep=0.11cm
    \begin{tabular}{|cc|c|}
    \hline
        Alpha range & Lightness range & F1-Score Top 3 \\ \hline
        (0.6, 0.8) & (0.6, 1.0) & \textbf{0.0908} \\ 
        (0.4, 0.6) & (0.2, 0.6) & 0.0965 \\ 
        (0.6, 0.8) & (0.4, 0.8) & 0.0075 \\ 
        (0.2, 0.4) & (0.6, 1.0)  & 0.0774 \\ 
        (0.4, 0.6) & (0.4, 0.8) & 0.0724 \\ 
        (0.2, 0.4) & (0.4, 0.8) & 0.0878 \\ 
        (0.2, 0.4) & (0.4, 0.8)  & 0.0624 \\ 
        (0.6, 0.8) & (0.2, 0.6)  & 0.0568 \\ 
        (0.2, 0.4) & (0.2, 0.6)  & 0.0849 \\ 
        (0.4, 0.6) & (0.6, 1.0)  & 0.0477 \\ \hline
    \end{tabular}
    
\end{table}

% \subsubsection{Motion Blur}
\begin{table}
\caption{Results of hyperparmeter tuning for Motion Blur.}
\label{tab:motion_blur_hp}
    \centering
    \small
    \tabcolsep=0.11cm
    \begin{tabular}{|c|c|}
    \hline
        Blur limit & F1-Score top-3 \\ \hline
        71 & 0.2689 \\ 
        61 & 0.2825 \\ 
        51 & 0.2707 \\ 
        41 & 0.2774 \\ 
        31 & 0.2831 \\ 
        21 & \textbf{0.3048} \\ 
        11 & 0.2872 \\ \hline
    \end{tabular}
    
\end{table}

% \subsubsection{Perspective}
\begin{table}
    \caption{Results of hyperparmeter tuning for Perspective.}
    \label{tab:perspective_hp}
    \centering
    \small
    \tabcolsep=0.11cm
    \begin{tabular}{|c|c|}
    \hline
    Scale & F1-Score top-3 \\ \hline
        0.7 & 0.2544 \\
        0.6 & 0.2542 \\
        0.5 & 0.2704 \\
        0.4 & 0.3014 \\
        0.3 & 0.3091 \\
        0.2 & \textbf{0.3126} \\ \hline
    \end{tabular}

\end{table}

\begin{table}[!htb]
\caption{Classification result when learning rate is $10^{-4}$ - a non-optimal classifier}    \label{tab:classification-result-2}
    \centering
    \small
    \tabcolsep=0.11cm
    \begin{tabular}{|l|lll|lll|}
    \hline
    & \multicolumn{3}{c}{\textbf{Top-3 Metrics}} & \multicolumn{3}{c}{\textbf{Top-1 Metrics}} \\ \hline
        \textbf{Method} & \textbf{\makecell{F1-score \\Top-3}} & \textbf{\makecell{Accuracy \\ Top-3}} & \textbf{\makecell{Precision \\ Top-3}} & \textbf{\makecell{F1-score \\ Top-1}} & \textbf{\makecell{Accuracy \\ Top-1}} & \textbf{\makecell{Precision \\ Top-1}} \\ \hline
        Baseline & 0.2784 & 0.4977 & 0.2264 & 0.3231 & 0.3434 & 0.3339 \\ \hline
        Rotation & \textbf{0.3364}& \textbf{0.5876} & \textbf{0.2585} & \textbf{0.432} & \textbf{0.4423} & \textbf{0.4586} \\ 
        Flipping & 0.3296 & 0.5826 & 0.2563 & 0.4104 & 0.4208 & 0.4351 \\ 
        Perspective & 0.286 & 0.5209 & 0.218 & 0.3339 & 0.3577 & 0.3414 \\ 
        CenterCrop 0.2 & 0.2829 & 0.5111 & 0.2204 & 0.3028 & 0.3252 & 0.3284 \\ 
        CenterCrop 0.4 & 0.2807 & 0.5205 & 0.209 & 0.3023 & 0.316 & 0.3333 \\ 
        MotionBlur & 0.3054 & 0.5209 & 0.2353 & 0.3233 & 0.3347 & 0.3486 \\ \hline
        CLAHE & 0.2569 & 0.4654 & 0.1971 & 0.2932 & 0.3027 & 0.3093 \\ 
        Sharpen & \textbf{0.3060} & 0.5095 & \textbf{0.2486} &  0.3168 & 0.3393 & 0.3525 \\ 
        ChannelShuffle & 0.2646 & 0.4682 & 0.2011 & 0.2945 & 0.314 & 0.3188 \\ 
        Solarize & 0.2251 & 0.4165 & 0.1666 & 0.2173 & 0.2354 & 0.2191 \\ 
        GaussNoise & 0.2868 & 0.5154 & 0.2168 & \textbf{0.3289} & \textbf{0.3498} & \textbf{0.3546} \\ 
        Normalize & 0.2778 & \textbf{0.5123} & 0.2152 & 0.3119 & 0.3315 & 0.3395 \\ 
        ToGray & 0.2638 & 0.4637 & 0.211 & 0.2787 & 0.2941 & 0.3011 \\ 
        ToSepia & 0.2695 & 0.4708 & 0.2111 & 0.3183 & 0.329 & 0.3436 \\ \hline
        Cutout & 0.2769 & 0.4891 & \textbf{0.2167} & \textbf{0.3199} & \textbf{0.3387} & 0.3396 \\ 
        GridDropout & \textbf{0.2827} & \textbf{0.4894} & \textbf{0.2167} & 0.3154 & 0.333 & \textbf{0.3454} \\ 
    \hline
    \end{tabular}
    
\end{table}